\clearpage
\section{Reconstructing balances at transition}
\label{sec:insurance-transition}

A disclosed restatement tells us what balances an insurer reported under the
new basis. It does not reveal the historical contract groups, assumptions or
services needed to reproduce their opening measurement. The three transition
methods address different limits on that reconstruction. The retained
December 2022 AASB~17 compilation provides a useful, explicitly scoped account
of those methods: its stated applicability covers annual periods beginning
from 1 January 2023 to before 1 July 2026. These provisions explain a design
requirement; they do not establish the rules governing Ping An's ASBE groups
\citep[Appendix C]{aasb17}.

First distinguish the dates. Initial application begins the annual period in
which the entity first applies the standard. Transition normally begins the
immediately preceding annual period. Thus a calendar-year entity first applying
the standard in 2023, with only 2022 presented as an adjusted comparative,
starts its transition calculation on 1 January 2022. If it elects to present
earlier adjusted comparatives, paragraph C25 instead takes transition back to
the beginning of the earliest such period. Neither date establishes when the
resulting report became public \citep[paragraphs C2 and C25--C27]{aasb17}.

Full retrospective application identifies and measures each group as if the
standard had always applied. Acquisition-cash-flow assets are considered
separately, balances that would not exist under the new rules are removed, and
the resulting net difference is recognised in equity. This route needs the
historical information required to reconstruct grouping, cash flows,
measurement assumptions and service already provided. It is required unless
application is impracticable or the specific paragraph C5A exception applies
\citep[paragraphs C3--C5]{aasb17}.

When full retrospection is impracticable for a group, the alternatives are
modified retrospection and fair-value transition. Modified retrospection aims
for the closest retrospective result possible using reasonable and supportable
information available without undue cost or effort. Each specified modification
is allowed only to the extent that the necessary retrospective information is
missing. If the information needed even for this modified method cannot be
obtained, paragraph C6(a) requires the fair-value approach. Missing history
therefore does not license an unrestricted choice of present assumptions
\citep[paragraphs C5--C8]{aasb17}.

For example, the modified route can reconstruct earlier expected cash flows
from estimates at transition, or an earlier determinable date, adjusted for
known intervening cash flows, including those of contracts that have expired.
Its discount-rate approximations require the specified yield-curve evidence;
its risk adjustment and allocation of past service require their own records.
The source also prescribes nil amounts for particular acquisition-cash-flow
items when the necessary information is unavailable. That narrow instruction
cannot be generalized into a rule that every missing amount is zero
\citep[paragraphs C12--C16]{aasb17}.

Fair-value transition takes a different starting point. Paragraph C20 determines
the CSM or loss component from the difference between the fair value of the
insurance group and its fulfilment cash flows, both measured at transition.
The insurance group's fair value is distinct from the value of its underlying
investment pool. Replacing one with the other would compute a different
quantity. The calculation needs a supported group valuation as well as the
cash-flow, discount and risk-adjustment inputs; C20 expressly excludes the
demand-feature rule in AASB~13 paragraph~47. Its permitted information-date,
grouping and accumulated-OCI choices also have specified conditions
\citep[paragraphs C20--C24B]{aasb17}.

\clearpage
There is one important qualification to the ordinary method choice. Under
C5A, a group with direct participation features may elect fair-value transition
even when retrospective application is possible, if the entity chooses the
B115 risk-mitigation option prospectively from transition and used qualifying
mitigation before that date. The applicable instrument, qualification and
designation conditions still matter. Treating every fair-value election as
proof of impracticability would therefore be wrong
\citep[paragraphs C3(b), C5A and B115--B118]{aasb17}.

The reconstruction must consequently retain the method and its basis for each
group, each modification used, the information date and evidence supporting
it, the separately measured FCF, CSM or loss component, relevant OCI and
acquisition balances, and the opening-equity reconciliation. Some transition
rules expressly permit transition-date information for specified assessments.
That permission concerns the accounting restatement; it does not make a later
report admissible in a forecast made before its publication. Detailed
impracticability judgments under AASB~108 and valuation under AASB~13, together
with actual group facts, remain necessary before this method description can
produce a supported opening measurement.

